\documentclass{article}

% ==============================================================================
% C++ Standard Paragraph & Sub-paragraph Numbering Macros
% ==============================================================================

\newcounter{stdpara}
\counterwithin*{stdpara}{section}
\counterwithin*{stdpara}{subsection}
\counterwithin*{stdpara}{subsubsection}

% Optional manual override: e.g., \setpnum{0} makes the next \pnum start at 1
\newcommand{\setparanum}[1]{\setcounter{stdpara}{#1}}

% Distance from the left text margin to the right edge of all margin numbers
\newlength{\marginnumsep}
\AtBeginDocument{\setlength{\marginnumsep}{1em}}

% Numbered paragraph command (\pnum)
\newcommand{\para}{%
  \par\medskip\noindent%
  \refstepcounter{stdpara}%
  \makebox[0pt][r]{\raisebox{0.35ex}{\scriptsize\thestdpara}\hspace{\marginnumsep}}%
  \ignorespaces
}

% Sub-paragraph counter and environment (subpara)
\newcounter{stdsubpara}[stdpara]
\renewcommand{\thestdsubpara}{\thestdpara.\arabic{stdsubpara}}

\newenvironment{subpara}{%
  \begin{list}{}{%
    \setlength{\topsep}{0.35em}%
    \setlength{\itemsep}{0.35em}%
    \setlength{\parsep}{0pt}%
    \settowidth{\labelwidth}{---}
    \setlength{\labelsep}{1em}
    \setlength{\itemindent}{1em}
    \setlength{\leftmargin}{\dimexpr\itemindent + \labelwidth + \labelsep\relax}%
    \setlength{\itemindent}{0pt}
    \renewcommand{\makelabel}[1]{
      \refstepcounter{stdsubpara}
      \makebox[0pt][r]{
        \footnotesize(\thestdsubpara)
        \hspace{\dimexpr\leftmargin - \labelsep - \labelwidth + \marginnumsep\relax}%
      }%
      ---%
    }%
  }%
}{%
  \end{list}%
}

% ==============================================================================
% Document Body
% ==============================================================================

\begin{document}

\section{General}
\subsection{Scope}

\para The set of \textit{diagnosable rules} consists of all syntactic and semantic rules in this document that are not explicitly referred to as ``undefined behavior''.

\para\label{p:impl-req} Although this document states only requirements on C++ implementations, those requirements are often easier to understand if they are phrased as requirements on programs. Such requirements have the following form:

\begin{subpara}
  \item\label{p:accept} If a program contains no violations of the rules in this document, a conforming implementation shall accept and correctly execute that program.
  \item If a program contains a violation of any diagnosable rule, a conforming implementation shall issue at least one diagnostic message.
  \item If a program contains a violation of a rule for which no diagnostic is required, this document places no requirement on implementations.
\end{subpara}

\para For classes and class templates, the library Clauses specify partial definitions (see paragraph~\ref{p:impl-req} and item~(\ref{p:accept})).

\subsection{Next Subsection}

\para Starting a new subsection automatically resets the paragraph counter back to 1.

\end{document}